% Document/LinearAlgebra/VectorSpaces/Matrices/MatrixOfMap.tex
% Phase 16e: Matrix of a Linear Map (The Commutative Diagram)

\newpage
\section{Matrix of a Linear Map}
\label{sec:matrix-of-map}

This section gives the matrix of an abstract linear map $T: U \to V$ with respect to
ordered bases. The commutative diagram is not the \emph{definition} of matrix
(that was Section~\ref{sec:mat-spaces}), but the \emph{definition} of ``the matrix of a
linear map'' --- and that is the concept that deserves emphasis.

\begin{remark}[{Why $K^{(I)}$?}]
    Abstract vector spaces have no canonical basis --- any basis requires an arbitrary choice.
    The spaces $K^{(I)}$ are the exception: they carry the canonical basis $(\CanonicalVector{i})_{i \in I}$,
    intrinsically defined by the function-space structure. This is the key reason we factor $T$
    through coordinate spaces: once we pass to $K^{(\kappa_1)} \to K^{(\kappa_2)}$, the
    canonical bases are automatically available and the Representation Theorem
    (Theorem~\ref{thm:rep-thm}) applies.
    \textbf{Matrices exist to exploit the canonical structure of $K^{(I)}$.}
\end{remark}

\subsection{Definition and Evaluation Formula}

Let $U$ and $V$ be $K$-vector spaces with ordered bases $\mathcal{B}_U = (b_j)_{j \in J}$ and
$\mathcal{B}_V = (c_i)_{i \in I}$, where $|J| = \kappa_1$ and $|I| = \kappa_2$.
Let $T: U \to V$ be $K$-linear. Recall the coordinate isomorphisms
$\Chart{\mathcal{B}_U}: U \xrightarrow{\sim} K^{(\kappa_1)}$ and
$\Chart{\mathcal{B}_V}: V \xrightarrow{\sim} K^{(\kappa_2)}$ from Theorem~\ref{thm:coord-bijection}.

\begin{definition}[Matrix of $T$ w.r.t.\ $\mathcal{B}_U$, $\mathcal{B}_V$]
\label{def:matrix-of-map}
    The \textbf{matrix of $T$ with respect to $\mathcal{B}_U$ and $\mathcal{B}_V$} is the
    $K$-linear map
    \[
        \CoordMatrix[\mathcal{B}_U][\mathcal{B}_V]{T}
        \defeq
        \Composition{\Chart{\mathcal{B}_V}}{\Composition{T}{\Frame{\mathcal{B}_U}}}
        \;\in\; \Hom[K]{K^{(\kappa_1)}}{K^{(\kappa_2)}}
    \]
    i.e.\ the composite
    $K^{(\kappa_1)} \xrightarrow{\Frame{\mathcal{B}_U}} U \xrightarrow{T} V
    \xrightarrow{\Chart{\mathcal{B}_V}} K^{(\kappa_2)}$.
    Since we identify $\Hom[K]{K^{(\kappa_1)}}{K^{(\kappa_2)}}$ with
    $\CFMat{\kappa_2}{\kappa_1}[K]$ via the Representation Theorem
    (Theorem~\ref{thm:rep-thm}), $\CoordMatrix[\mathcal{B}_U][\mathcal{B}_V]{T}$ is
    simultaneously a linear map between canonical spaces and a column-finite matrix.
    The commutative square is:
    \[
    \begin{tikzcd}
        U \arrow[r, "T"] \arrow[d, "\Chart{\mathcal{B}_U}"'] &
        V \arrow[d, "\Chart{\mathcal{B}_V}"] \\
        K^{(\kappa_1)} \arrow[r,
            "{\CoordMatrix[\mathcal{B}_U][\mathcal{B}_V]{T}}"'] &
        K^{(\kappa_2)}
    \end{tikzcd}
    \]
\end{definition}

\begin{remark}
    \leavevmode
    \begin{itemize}
        \item Column-finiteness is automatic:
              $\CoordMatrix[\mathcal{B}_U][\mathcal{B}_V]{T}$ maps into $K^{(\kappa_2)}$
              (since $\Chart{\mathcal{B}_V}$ does), so the Representation Theorem places it
              in $\CFMat{\kappa_2}{\kappa_1}[K]$.
        \item This definition requires \textbf{ordered} bases (to have a specific coordinate
              isomorphism); unordered bases give only an isomorphism up to reordering.
        \item \textbf{Notation}: Always write
              $\CoordMatrix[\mathcal{B}_U][\mathcal{B}_V]{T}$ with basis indices, even when the
              bases are canonical:
              $\CoordMatrix[\CanonicalBase{\kappa_1}][\CanonicalBase{\kappa_2}]{T}$.
    \end{itemize}
\end{remark}

\begin{theorem}[Evaluation Formula]
\label{thm:column-formula}
    The $j$-th column of $\CoordMatrix[\mathcal{B}_U][\mathcal{B}_V]{T}$ (i.e., the value
    of the associated linear map on $\CanonicalVector{j} \in K^{(\kappa_1)}$) equals the coordinate vector
    $\Chart{\mathcal{B}_V}(T(b_j))$. Explicitly, if
    $T(b_j) = \sum_{i \in I} a_{ij}\, c_i$, then the $(i,j)$-entry of the matrix is $a_{ij}$.
\end{theorem}

\begin{proof}
    The composite evaluated at $\CanonicalVector{j}$ gives
    $\Chart{\mathcal{B}_V}(T(\Frame{\mathcal{B}_U}(\CanonicalVector{j})))
    = \Chart{\mathcal{B}_V}(T(b_j))$.
    The $i$-th component is
    $(\Chart{\mathcal{B}_V}(T(b_j)))_i = a_{ij}$.
\end{proof}

\begin{remark}[Practical Consequence]
    To compute $\CoordMatrix[\mathcal{B}_U][\mathcal{B}_V]{T}$: apply $T$ to each basis vector
    $b_j \in \mathcal{B}_U$, express the result $T(b_j)$ in the basis $\mathcal{B}_V$,
    and place those coordinates as the $j$-th column:
    \[
        \CoordMatrix[\mathcal{B}_U][\mathcal{B}_V]{T}
        = \begin{pmatrix}
            \vert & \vert & & \vert \\[4pt]
            \CoordVec[\mathcal{B}_V]{T(b_1)} & \CoordVec[\mathcal{B}_V]{T(b_2)}
            & \cdots & \CoordVec[\mathcal{B}_V]{T(b_{\kappa_1})} \\[4pt]
            \vert & \vert & & \vert
        \end{pmatrix}
    \]
    where $\mathcal{B}_U = (b_1, b_2, \ldots, b_{\kappa_1})$.
\end{remark}

\subsection{Coordinate Formula and Composition}

\begin{corollary}[Coordinate Formula]
\label{prop:coord-formula}
    For any $v \in U$:
    \[
        \CoordVec[\mathcal{B}_V]{T(v)}
        = \CoordMatrix[\mathcal{B}_U][\mathcal{B}_V]{T}
          \cdot \CoordVec[\mathcal{B}_U]{v}
    \]
\end{corollary}

\begin{proof}
    This is the matrix-times-column formula of Theorem~\ref{thm:rep-thm} applied to the
    linear map
    $\Composition{\Chart{\mathcal{B}_V}}{\Composition{T}{\Frame{\mathcal{B}_U}}}$:
    \[
        \CoordMatrix[\mathcal{B}_U][\mathcal{B}_V]{T} \cdot \Chart{\mathcal{B}_U}(v)
        = (\Composition{\Chart{\mathcal{B}_V}}{\Composition{T}{\Frame{\mathcal{B}_U}}})
          (\Chart{\mathcal{B}_U}(v))
        = \Chart{\mathcal{B}_V}(T(v))
        = \CoordVec[\mathcal{B}_V]{T(v)} \qedhere
    \]
\end{proof}

\begin{remark}[Matrix-Times-Vector]
    The formula
    $\CoordMatrix[\mathcal{B}_U][\mathcal{B}_V]{T} \cdot x = Ax$ is a \textbf{theorem},
    not a definition. It follows from the Representation Theorem and linearity.
    Explicitly, if $A_{ij}$ are the entries and $x = \CoordVec[\mathcal{B}_U]{v}$, then
    \[
        \bigl(\CoordVec[\mathcal{B}_V]{T(v)}\bigr)_i = \sum_{j} A_{ij}\, x_j
    \]
    This is the standard row-by-column product $(Ax)_i = \sum_j A_{ij} x_j$.
\end{remark}

\begin{corollary}[Composition Equals Multiplication]
\label{thm:comp-mult}
    For composable $T: U \to V$, $S: V \to W$ with ordered bases
    $\mathcal{B}_U$, $\mathcal{B}_V$, $\mathcal{B}_W$:
    \[
        \CoordMatrix[\mathcal{B}_U][\mathcal{B}_W]{\Composition{S}{T}}
        = \Composition{\CoordMatrix[\mathcal{B}_V][\mathcal{B}_W]{S}}
          {\CoordMatrix[\mathcal{B}_U][\mathcal{B}_V]{T}}
    \]
    That is, composition of abstract maps becomes composition (equivalently, matrix
    multiplication) of their coordinate representations.
\end{corollary}

\begin{proof}
    By definition, inserting
    $\Identity{V} = \Composition{\Frame{\mathcal{B}_V}}{\Chart{\mathcal{B}_V}}$:
    \begin{align*}
        \CoordMatrix[\mathcal{B}_U][\mathcal{B}_W]{\Composition{S}{T}}
        &= \Composition{\Chart{\mathcal{B}_W}}
           {\Composition{(\Composition{S}{T})}{\Frame{\mathcal{B}_U}}} \\
        &= \Composition{\Chart{\mathcal{B}_W}}
           {\Composition{S}
           {\Composition{\Frame{\mathcal{B}_V}}
           {\Composition{\Chart{\mathcal{B}_V}}}
           {\Composition{T}{\Frame{\mathcal{B}_U}}}}} \\
        &= \Composition{(\Composition{\Chart{\mathcal{B}_W}}
           {\Composition{S}{\Frame{\mathcal{B}_V}}})}
           {(\Composition{\Chart{\mathcal{B}_V}}
           {\Composition{T}{\Frame{\mathcal{B}_U}}})} \\
        &= \Composition{\CoordMatrix[\mathcal{B}_V][\mathcal{B}_W]{S}}
           {\CoordMatrix[\mathcal{B}_U][\mathcal{B}_V]{T}} \qedhere
    \end{align*}
\end{proof}

\subsection{Vectors as Maps}

\begin{remark}[Matrix-Times-Vector as Composition via {$V \cong \Hom[K]{K}{V}$}]
\label{rem:mat-times-vec-composition}
    The identification $V \cong \Hom[K]{K}{V}$ from
    Remark~\ref{rem:vec-as-hom} explains the coordinate formula as a special case
    of composition equals multiplication. For $T: U \to V$ and $u \in U$, the composition
    $\Composition{T}{\lambda_u}: K \to V$ sends $1 \mapsto T(u)$, so
    $\CoordMatrix[(1)][\mathcal{B}_V]{\Composition{T}{\lambda_u}}
    = \CoordVec[\mathcal{B}_V]{T(u)}$.
    By Corollary~\ref{thm:comp-mult}:
    \[
        \CoordVec[\mathcal{B}_V]{T(u)}
        = \CoordMatrix[\mathcal{B}_U][\mathcal{B}_V]{T}
          \cdot \CoordMatrix[(1)][\mathcal{B}_U]{\lambda_u}
        = \CoordMatrix[\mathcal{B}_U][\mathcal{B}_V]{T}
          \cdot \CoordVec[\mathcal{B}_U]{u}
    \]
    The shared summation index --- column index of
    $\CoordMatrix[\mathcal{B}_U][\mathcal{B}_V]{T}$ paired with row index of
    $\CoordVec[\mathcal{B}_U]{u}$ --- is forced by composition:
    $\Hom[K]{U}{V} \times \Hom[K]{K}{U} \to \Hom[K]{K}{V}$ contracts the middle space $U$.
    \textbf{Matrix-times-vector is not a separate operation; it is matrix multiplication
    applied to a column matrix.}
\end{remark}

\begin{remark}[Coordinate Vectors and Matrices as the Same Construction]
    Both coordinate vectors and matrices are instances of the same construction:
    \textbf{transport to canonical spaces via coordinate isomorphisms}.
    \begin{itemize}
        \item For a vector $v \in V$ (equivalently, a map
              $\lambda_v \in \Hom[K]{K}{V}$ via Theorem~\ref{thm:hom-k-v-iso}),
              the chain is
              \[
                  V \xrightarrow{\;\Chart{\mathcal{B}}\;} K^{(I)}
                  \;\cong\; \CFMat{I}{\{*\}}[K]
              \]
              The coordinate vector $\CoordVec[\mathcal{B}]{v} = \Chart{\mathcal{B}}(v)$
              is a column matrix with row-index set $I$ and a singleton column-index set.
        \item For a map $T \in \Hom[K]{U}{V}$, the chain is
              \[
                  \Hom[K]{U}{V}
                  \xrightarrow{\;\Chart{\mathcal{B}_V} \circ (-) \circ
                  \Frame{\mathcal{B}_U}\;}
                  \Hom[K]{K^{(\kappa_1)}}{K^{(\kappa_2)}}
                  \;\cong\; \CFMat{\kappa_2}{\kappa_1}[K]
              \]
              The matrix
              $\CoordMatrix[\mathcal{B}_U][\mathcal{B}_V]{T}$
              is the image of $T$ under this composite.
    \end{itemize}
    Both chains have the same structure: abstract object $\to$ concrete canonical
    space $\to$ grid of numbers. We usually \emph{stop one step earlier} for vectors,
    because a column vector already looks like a (degenerate) grid --- there is nothing to
    gain by writing it as a $\kappa \times 1$ matrix.

    The apparent asymmetry in the notation reflects only the degeneracy of the source basis
    when the source is $K$: the notation $\CoordVec[\mathcal{B}]{v}$ carries one basis
    superscript because the source basis $\{1\}$ of $K$ is canonical and requires no choice.
    The notation $\CoordMatrix[\mathcal{B}_U][\mathcal{B}_V]{T}$ carries two because both
    source and target have non-trivial bases.
\end{remark}

\subsection{General Hom via Bases}

\begin{theorem}[{$\Hom[K]{U}{V} \cong \CFMat{|\mathcal{B}_V|}{|\mathcal{B}_U|}[K]$}]
\label{thm:hom-cfmat}
    For $U$, $V$ with ordered bases $\mathcal{B}_U$, $\mathcal{B}_V$, the map
    $T \mapsto \CoordMatrix[\mathcal{B}_U][\mathcal{B}_V]{T}$ is a $K$-linear isomorphism:
    \[
        \Hom[K]{U}{V} \xrightarrow{\;\sim\;}
        \CFMat{|\mathcal{B}_V|}{|\mathcal{B}_U|}[K]
    \]
    The isomorphism lands in the \textbf{column-finite} space: for any
    $T \in \Hom[K]{U}{V}$ and any $b_j \in \mathcal{B}_U$, the image $T(b_j) \in V$ has
    finite support in $\mathcal{B}_V$ (by definition of a basis), so the $j$-th column of
    $\CoordMatrix[\mathcal{B}_U][\mathcal{B}_V]{T}$ is finitely supported.
\end{theorem}

\begin{proof}
    Compose the coordinate isomorphisms:
    \[
    \begin{tikzcd}[column sep=huge]
        \Hom[K]{U}{V}
        \arrow[r, "\sim", "{\Chart{\mathcal{B}_V} \circ (-) \circ \Frame{\mathcal{B}_U}}"']
        & \Hom[K]{K^{(|\mathcal{B}_U|)}}{K^{(|\mathcal{B}_V|)}}
        \arrow[r, "\sim", "\Phi"']
        & \CFMat{|\mathcal{B}_V|}{|\mathcal{B}_U|}[K]
    \end{tikzcd}
    \]
    The first map sends $T$ to
    $\Composition{\Chart{\mathcal{B}_V}}{\Composition{T}{\Frame{\mathcal{B}_U}}}$
    and is a $K$-linear isomorphism (conjugation by fixed isomorphisms). The second is $\Phi$
    from Theorem~\ref{thm:rep-thm}. Their composition is the desired isomorphism.
\end{proof}

\begin{remark}
    Different choices of bases yield different but isomorphic identifications (related by
    change-of-basis matrices).
\end{remark}

\begin{remark}[{Coordinate Vectors as Matrix Representations --- Deferred from
    Remark~\ref{rem:vec-as-hom}}]
\label{rem:coord-vec-as-matrix}
    The Representation Theorem (Theorem~\ref{thm:rep-thm}) showed that
    $\CFMat{\kappa_2}{1}[K] \cong K^{(\kappa_2)}$: column matrices \emph{are} vectors. Now
    that we have the general notation $\CoordMatrix[\mathcal{B}_U][\mathcal{B}_V]{T}$, we
    can state the full picture: for any $V$ with ordered basis
    $\mathcal{B}_V = (c_i)_{i \in I}$, the coordinate vector
    $\CoordVec[\mathcal{B}_V]{v}$ \textbf{is} the matrix
    $\CoordMatrix[(1)][\mathcal{B}_V]{\lambda_v}$ --- an $I \times \{1\}$ column matrix,
    where $\lambda_v: K \to V$ sends $1 \mapsto v$. The coordinate representation of a
    vector and the matrix representation of a linear map are literally the same object.
\end{remark}

\begin{remark}[Three Layers of a Linear Map]\label{rem:three-layers-map}
    The matrix of a linear map exhibits the same three-layer structure as the coordinate
    vector of a vector (Remark~\ref{rem:three-layers-vector}):
    \begin{enumerate}
        \item \textbf{The geometric object}: $T: U \to V$, a $K$-linear map between abstract
              vector spaces. Basis-free.
        \item \textbf{The coordinate representation}:
              $\Composition{\Chart{\mathcal{B}_V}}{\Composition{T}{\Frame{\mathcal{B}_U}}}
              \in \Hom[K]{K^{(\kappa_1)}}{K^{(\kappa_2)}}$, a $K$-linear map between canonical
              spaces. This is an algebraic object: it can be composed, it has a kernel and
              image, and the structural theorems (composition = multiplication, change of basis)
              are proved at this level.
        \item \textbf{The data structure}:
              $\CoordMatrix[\mathcal{B}_U][\mathcal{B}_V]{T}
              \in \CFMat{\kappa_2}{\kappa_1}[K]$, the grid of scalars extracted via the
              Representation Theorem ($\Phi$). This is the object we compute with: we write it
              on a page, store it in a computer, perform Gaussian elimination on it, compute
              its determinant.
    \end{enumerate}
    Layers (1) and (2) are related by the coordinate isomorphisms ($\Chart{\mathcal{B}}$).
    Layers (2) and (3) are the same set with different structure, related by the Representation
    Theorem ($\Phi$). We proved the structural theorems --- composition, change of basis, rank
    invariance --- at layer (2), where they are tautologies about function composition. The
    entry formulas are then \emph{consequences} extracted at layer (3).

    For vectors, we noted that layers (2) and (3) are the same set ($K^{(I)}$) viewed as a
    vector space or as a set of tuples. For linear maps, the situation is identical:
    $\Hom[K]{K^{(\kappa_1)}}{K^{(\kappa_2)}}$ and $\CFMat{\kappa_2}{\kappa_1}[K]$ are the
    same set viewed as a $K$-vector space (with composition) or as a set of grids (with
    entrywise operations). Having established the identification, we pass freely between the
    two views.
\end{remark}

\begin{example}[Matrix of a Linear Map --- Standard Bases]
\label{ex:matrix-of-map}
    Let $T: \bb{R}^3 \to \bb{R}^2$ be defined by $T(x,y,z) = (x + 2y, 3z)$. With the
    standard ordered bases $\mathcal{B}_U = (\CanonicalVector{0}, \CanonicalVector{1}, \CanonicalVector{2})$ on $\bb{R}^3$ and
    $\mathcal{B}_V = (\CanonicalVector{0}, \CanonicalVector{1})$ on $\bb{R}^2$:
    \begin{align*}
        T(\CanonicalVector{0}) &= (1, 0) = 1 \cdot \CanonicalVector{0} + 0 \cdot \CanonicalVector{1} \\
        T(\CanonicalVector{1}) &= (2, 0) = 2 \cdot \CanonicalVector{0} + 0 \cdot \CanonicalVector{1} \\
        T(\CanonicalVector{2}) &= (0, 3) = 0 \cdot \CanonicalVector{0} + 3 \cdot \CanonicalVector{1}
    \end{align*}
    Placing these as columns:
    \[
        \CoordMatrix[\mathcal{B}_U][\mathcal{B}_V]{T}
        = \begin{pmatrix} 1 & 2 & 0 \\ 0 & 0 & 3 \end{pmatrix}
    \]
\end{example}

\begin{example}[Matrix of a Linear Map --- Non-Standard Basis]
\label{ex:matrix-of-map-nonstandard}
    Let $D: \bb{R}_{\leq 2}[x] \to \bb{R}_{\leq 2}[x]$ be the derivative on polynomials of
    degree $\leq 2$, with ordered basis $\mathcal{B} = (1, x, x^2)$:
    \begin{align*}
        D(1) &= 0 = 0 \cdot 1 + 0 \cdot x + 0 \cdot x^2 \\
        D(x) &= 1 = 1 \cdot 1 + 0 \cdot x + 0 \cdot x^2 \\
        D(x^2) &= 2x = 0 \cdot 1 + 2 \cdot x + 0 \cdot x^2
    \end{align*}
    So:
    \[
        \CoordMatrix[\mathcal{B}][\mathcal{B}]{D}
        = \begin{pmatrix} 0 & 1 & 0 \\ 0 & 0 & 2 \\ 0 & 0 & 0 \end{pmatrix}
    \]
\end{example}

\subsection{Examples of the Coordinate Formula}

\begin{example}[{$\CoordVec[\mathcal{B}_V]{T(v)} = \CoordMatrix[\mathcal{B}_U][\mathcal{B}_V]{T} \cdot \CoordVec[\mathcal{B}_U]{v}$} --- I]
    With $T$, $\mathcal{B}_U$, $\mathcal{B}_V$ as in Example~\ref{ex:matrix-of-map}, let
    $v = (1, -1, 2) \in \bb{R}^3$. Then $\CoordVec[\mathcal{B}_U]{v}
    = \begin{psmallmatrix} 1 \\ -1 \\ 2 \end{psmallmatrix}$ and:
    \[
        \CoordMatrix[\mathcal{B}_U][\mathcal{B}_V]{T} \cdot \CoordVec[\mathcal{B}_U]{v}
        = \begin{pmatrix} 1 & 2 & 0 \\ 0 & 0 & 3 \end{pmatrix}
          \begin{pmatrix} 1 \\ -1 \\ 2 \end{pmatrix}
        = \begin{pmatrix} 1 - 2 \\ 6 \end{pmatrix}
        = \begin{pmatrix} -1 \\ 6 \end{pmatrix}
    \]
    Direct check: $T(1, -1, 2) = (1 + 2(-1),\, 3 \cdot 2) = (-1, 6)$, so
    $\CoordVec[\mathcal{B}_V]{T(v)}
    = \begin{psmallmatrix} -1 \\ 6 \end{psmallmatrix}$. $\checkmark$
\end{example}

\begin{example}[{$\CoordVec[\mathcal{B}]{D(p)} = \CoordMatrix[\mathcal{B}][\mathcal{B}]{D} \cdot \CoordVec[\mathcal{B}]{p}$} --- II]
    With $D$ and $\mathcal{B} = (1, x, x^2)$ as in
    Example~\ref{ex:matrix-of-map-nonstandard}, let $p = 3 + x - 2x^2$. Then
    $\CoordVec[\mathcal{B}]{p}
    = \begin{psmallmatrix} 3 \\ 1 \\ -2 \end{psmallmatrix}$ and:
    \[
        \CoordMatrix[\mathcal{B}][\mathcal{B}]{D} \cdot \CoordVec[\mathcal{B}]{p}
        = \begin{pmatrix} 0 & 1 & 0 \\ 0 & 0 & 2 \\ 0 & 0 & 0 \end{pmatrix}
          \begin{pmatrix} 3 \\ 1 \\ -2 \end{pmatrix}
        = \begin{pmatrix} 1 \\ -4 \\ 0 \end{pmatrix}
    \]
    Direct check: $D(3 + x - 2x^2) = 1 - 4x$, so
    $\CoordVec[\mathcal{B}]{D(p)}
    = \begin{psmallmatrix} 1 \\ -4 \\ 0 \end{psmallmatrix}$. $\checkmark$
\end{example}

\subsection{Examples of Composition Equals Multiplication}

\begin{example}[{$\CoordMatrix[\mathcal{B}_U][\mathcal{B}_W]{\Composition{S}{T}} = \CoordMatrix[\mathcal{B}_V][\mathcal{B}_W]{S} \cdot \CoordMatrix[\mathcal{B}_U][\mathcal{B}_V]{T}$} --- I]
    Let $T: \bb{R}^2 \to \bb{R}^2$ be $T(x,y) = (x + y, x - y)$ and
    $S: \bb{R}^2 \to \bb{R}^2$ be $S(x,y) = (2x, y)$, all with the standard basis
    $\mathcal{B} = (\CanonicalVector{0}, \CanonicalVector{1})$.
    \[
        \CoordMatrix[\mathcal{B}][\mathcal{B}]{T}
        = \begin{pmatrix} 1 & 1 \\ 1 & -1 \end{pmatrix}, \qquad
        \CoordMatrix[\mathcal{B}][\mathcal{B}]{S}
        = \begin{pmatrix} 2 & 0 \\ 0 & 1 \end{pmatrix}
    \]
    Matrix product:
    \[
        \CoordMatrix[\mathcal{B}][\mathcal{B}]{S}
        \cdot \CoordMatrix[\mathcal{B}][\mathcal{B}]{T}
        = \begin{pmatrix} 2 & 0 \\ 0 & 1 \end{pmatrix}
          \begin{pmatrix} 1 & 1 \\ 1 & -1 \end{pmatrix}
        = \begin{pmatrix} 2 & 2 \\ 1 & -1 \end{pmatrix}
    \]
    Direct check: $(\Composition{S}{T})(x,y) = S(x+y, x-y) = (2(x+y), x-y)
    = (2x + 2y, x - y)$, so
    \[
        \CoordMatrix[\mathcal{B}][\mathcal{B}]{\Composition{S}{T}}
        = \begin{pmatrix} 2 & 2 \\ 1 & -1 \end{pmatrix} \quad \checkmark
    \]
\end{example}

\begin{example}[{$\CoordMatrix[\mathcal{B}][\mathcal{B}]{\Composition{D}{D}} = \bigl(\CoordMatrix[\mathcal{B}][\mathcal{B}]{D}\bigr)^2$} --- II]
    With $D$ and $\mathcal{B} = (1, x, x^2)$ as in
    Example~\ref{ex:matrix-of-map-nonstandard}:
    \[
        \bigl(\CoordMatrix[\mathcal{B}][\mathcal{B}]{D}\bigr)^2
        = \begin{pmatrix} 0 & 1 & 0 \\ 0 & 0 & 2 \\ 0 & 0 & 0 \end{pmatrix}
          \begin{pmatrix} 0 & 1 & 0 \\ 0 & 0 & 2 \\ 0 & 0 & 0 \end{pmatrix}
        = \begin{pmatrix} 0 & 0 & 2 \\ 0 & 0 & 0 \\ 0 & 0 & 0 \end{pmatrix}
    \]
    Direct check: $D^2(1) = 0$, $D^2(x) = D(1) = 0$, $D^2(x^2) = D(2x) = 2$, so
    \[
        \CoordMatrix[\mathcal{B}][\mathcal{B}]{D^2}
        = \begin{pmatrix} 0 & 0 & 2 \\ 0 & 0 & 0 \\ 0 & 0 & 0 \end{pmatrix}
        \quad \checkmark
    \]
\end{example}
